\clearpage
\SourceLeaf{615}
\ArithmeticHead{586}{L}{INHALT.}
\begingroup
\newcommand{\ContentsRow}[5]{\par\noindent\hangindent=#1\hspace*{#1}#2\hspace{.4em}\mbox{\footnotesize #3}\hfill\mbox{#4\quad #5}\par}
\ContentsRow{0pt}{109.\quad Mehrdeutigkeit des agM., 1800 Juni 3}{.\,.\,.\,.\,.\,.\,.\,.\,.\,.\,.\,.\,.\,.\,.\,.\,.\,.\,.}{Seite}{550}
\ContentsRow{0pt}{110.\quad Anwendung auf die elliptischen Transzendenten, 1800 Juni 5}{.\,.\,.\,.\,.\,.\,.\,.\,.\,.\,.}{—}{551}
\ContentsRow{0pt}{111.\quad Rektifikation der Ellipse, 1800 Juni 10}{.\,.\,.\,.\,.\,.\,.\,.\,.\,.\,.\,.\,.\,.\,.\,.\,.\,.\,.}{—}{551}
\ContentsRow{0pt}{112.\quad Numerische Berechnung von Exponentialgrößen, 1800 Juni 12}{.\,.\,.\,.\,.\,.\,.\,.\,.\,.}{—}{551}
\ContentsRow{0pt}{113.\quad Eine Aufgabe der Wahrscheinlichkeitsrechnung, 1800 Oktober 25}{.\,.\,.\,.\,.\,.\,.\,.\,.}{—}{552}
\ContentsRow{0pt}{114, 115.\quad Bestimmung der Klassenanzahl binärer quadratischer Formen, 1800 November 30, Dezember 3}{.\,.\,.\,.\,.\,.\,.\,.\,.\,.\,.\,.\,.\,.\,.\,.\,.\,.\,.\,.\,.\,.\,.\,.\,.\,.\,.\,.}{—}{556}
\ContentsRow{0pt}{116.\quad Unvermeidlichkeit der Hilfsgleichungen bei der Kreisteilung, 1801 April 6}{.\,.\,.\,.\,.\,.}{—}{556}
\ContentsRow{0pt}{117.\quad Formel für das jüdische Osterfest, 1801 April 1}{.\,.\,.\,.\,.\,.\,.\,.\,.\,.\,.\,.\,.\,.\,.\,.}{—}{560}
\ContentsRow{0pt}{118.\quad Zum fünften Beweise des Reziprozitätssatzes der quadratischen Reste, 1801 Mitte Mai}{.\,.}{—}{560}
\ContentsRow{0pt}{119.\quad Bahnbestimmung der Himmelskörper, 1801 Mitte September}{.\,.\,.\,.\,.\,.\,.\,.\,.\,.\,.}{—}{561}
\ContentsRow{0pt}{120.\quad Mondbewegung, 1801 August}{.\,.\,.\,.\,.\,.\,.\,.\,.\,.\,.\,.\,.\,.\,.\,.\,.\,.\,.\,.\,.\,.}{—}{562}
\ContentsRow{0pt}{121.\quad Neues zur theoretischen Astronomie, 1801 Oktober}{.\,.\,.\,.\,.\,.\,.\,.\,.\,.\,.\,.\,.\,.\,.}{—}{563}
\ContentsRow{0pt}{122.\quad Astronomische Tätigkeit während der Jahre 1802, 1803, 1804}{.\,.\,.\,.\,.\,.\,.\,.\,.\,.\,.}{—}{563}
\ContentsRow{0pt}{123.\quad Fünfter Beweis des Reziprozitätssatzes der quadratischen Reste, 1805 August 30}{.\,.\,.\,.}{—}{563}
\ContentsRow{0pt}{124.\quad Über Interpolation, 1805 November}{.\,.\,.\,.\,.\,.\,.\,.\,.\,.\,.\,.\,.\,.\,.\,.\,.\,.\,.\,.}{—}{564}
\ContentsRow{0pt}{125, 126, 127.\quad Bahnbestimmung der Himmelskörper, 1806 Januar bis Mai}{.\,.\,.\,.\,.\,.\,.\,.}{—}{564,565}
\ContentsRow{0pt}{128.\quad Biquadratische Resolvente der Kreisteilungsgleichung, 1806 April bis Mai}{.\,.\,.\,.\,.\,.}{—}{565}
\ContentsRow{0pt}{129.\quad Bahnbestimmung eines Planeten aus vier geocentrischen Orten, 1807 Januar 21}{.\,.\,.\,.}{—}{565}
\ContentsRow{0pt}{130—133.\quad Biquadratische und kubische Reste, 1807 Februar 15, 17, 22, 24}{.\,.\,.\,.\,.\,.\,.}{—}{565,566}
\ContentsRow{0pt}{134.\quad Sechster Beweis des Reziprozitätssatzes der quadratischen Reste, 1807 Mai 6}{.\,.\,.\,.\,.}{—}{566}
\ContentsRow{0pt}{135.\quad Die kubische Resolvente der Kreisteilungsgleichung, 1808 Mai 10}{.\,.\,.\,.\,.\,.\,.\,.\,.}{—}{567}
\ContentsRow{0pt}{136.\quad Irreduzibilität der Kreisteilungsgleichung für einen zusammengesetzten Grad, 1808 Juni 12}{}{—}{567}
\ContentsRow{0pt}{137.\quad Kubische Formen, 1808 Dezember 23}{.\,.\,.\,.\,.\,.\,.\,.\,.\,.\,.\,.\,.\,.\,.\,.\,.\,.\,.}{—}{568}
\ContentsRow{0pt}{138.\quad Kubische Reste, 1809 Januar 6}{.\,.\,.\,.\,.\,.\,.\,.\,.\,.\,.\,.\,.\,.\,.\,.\,.\,.\,.\,.\,.}{—}{568}
\ContentsRow{0pt}{139.\quad Reihen zum agM., 1809 Juni 20}{.\,.\,.\,.\,.\,.\,.\,.\,.\,.\,.\,.\,.\,.\,.\,.\,.\,.\,.\,.\,.}{—}{569}
\ContentsRow{0pt}{140.\quad Fünfteilung des agM., 1809 Juni 29}{.\,.\,.\,.\,.\,.\,.\,.\,.\,.\,.\,.\,.\,.\,.\,.\,.\,.\,.\,.}{—}{569}
\ContentsRow{0pt}{141.\quad Über den zweiten Beweis der Wurzelexistenz, 1812 Februar 29}{.\,.\,.\,.\,.\,.\,.\,.\,.\,.}{—}{569}
\ContentsRow{0pt}{142, 143.\quad Anziehung der elliptischen Sphäroide, 1812 September 26, Oktober 15}{.\,.\,.\,.\,.\,.}{—}{570}
\ContentsRow{0pt}{144, 145.\quad Biquadratische Reste; Kometenbahnen, 1813 Oktober 23}{.\,.\,.\,.\,.\,.\,.\,.\,.\,.\,.}{—}{571}
\ContentsRow{0pt}{146.\quad Zusammenhang zwischen biquadratischen Resten und lemniskatischen Funktionen, 1814 Juli 9}{.\,.\,.\,.\,.\,.\,.\,.\,.\,,\,.\,.\,.\,.\,.\,.\,.\,.\,.\,.\,.\,.\,.\,.\,.\,.\,.\,.\,.}{—}{571}
\ContentsRow{0pt}{\emph{Schlussbemerkung}}{.\,.\,.\,.\,.\,.\,.\,.\,.\,.\,.\,.\,.\,.\,.\,.\,.\,.\,.\,.\,.\,.\,.\,.\,.\,.\,.\,.\,.}{—}{572}
\ContentsRow{0pt}{\emph{Sachverzeichnis zum Tagebuch}}{.\,.\,.\,.\,.\,.\,.\,.\,.\,.\,.\,.\,.\,.\,.\,.\,.\,.\,.\,.\,.\,.\,.\,.}{—}{573}
\begin{center}\rule{.13\linewidth}{.4pt}\end{center}
\ContentsRow{0pt}{\emph{Berichtigungen}}{.\,.\,.\,.\,.\,.\,.\,.\,.\,.\,.\,.\,.\,.\,.\,.\,.\,.\,.\,.\,.\,.\,.\,.\,.\,.\,.\,.\,.\,.}{—}{575}
\ContentsRow{0pt}{\emph{Bemerkungen zur ersten Abteilung des zehnten Bandes}}{.\,.\,.\,.\,.\,.\,.\,.\,.\,.\,.\,.\,.\,.\,.}{—}{576}
\begin{center}\rule{.26\linewidth}{.4pt}\end{center}
\par\bigskip\begin{center}\small Göttingen, Druck der Dieterichschen Univ.-Buchdruckerei (W. Fr. Kaestner).\end{center}
\endgroup
